\section{Conclusion}

This paper demonstrated the effectiveness of visual instruction tuning.
We presented an automatic pipeline to create language-image instruction-following data, based on which we train \shortname{}, a multimodal model to follow human intent to complete visual tasks. It achieves the new SoTA accuracy when fine-tuned on ScienceQA, and excellent visual chat capabilities when fine-tuned on multimodal chat data.
Besides, we present the first benchmark to study multimodal instruction-following capability.
This paper is an initial step in visual instruction tuning, and mainly focuses on real-life tasks. For more quantitative results of \shortname{} on academic benchmarks, please refer to the improved baselines with visual instruction tuning~\cite{liu2023improvedllava}. We hope our work can inspire future research on building more capable multimodal models.

\textbf{Acknowledgements.} 
We thank Baolin Peng and Pan Lu for valuable discussions on instruction-tuning language models and Science QA, respectively. 
We thank the LLaMA team for giving us access to their models, and open-source projects, including Alpaca and Vicuna.
This work was supported in part by NSF CAREER IIS2150012, and Institute of Information \& communications Technology Planning \& Evaluation(IITP) grants funded by the Korea government(MSIT) (No. 2022-0-00871, Development of AI Autonomy and Knowledge Enhancement for AI Agent Collaboration) and (No. RS-2022-00187238, Development of Large Korean Language Model Technology for Efficient Pre-training).